\scanpage{8}
% archive foliation: 7
\uline{Prop} Soit $L_*$ un $k$-module $\frac12$ simplicial.
Les conditions suivantes sont équivalentes

a) les $L_n$ et $\pi_n(L_*)$ sont des $k$-modules
projectifs.

a') Les composantes \deleted{\unsure{de}} $(L^!_*)_n$ de $L^!_*$ (le
complexe de chaînes de Dold-Puppe associé à
$L_*$) et les $H_n(L^!_*)$ sont des $k$-modules
projectifs

b) $L_*$ est isomorphe à un \deleted{\unsure{somme}} \inserted{produit}
direct \deleted{de modules} de la forme
\[
\left(\prod_{i\geq1}\Lambda^i\Phi_{*k}\otimes_k M_i\right)
\times\prod_{j\geq0}(\Lambda^j\Psi_{*k})\otimes_k N_j
\]
\unsure{et} b') $L^!$ est isomorphe à un produit
direct (de complexes de cochaînes) de
la forme
\diagram{\displaystyle
\prod_{i\geq1}\left(
\begin{tikzcd}[column sep=small]
0\arrow[r]&\underset{\text{en degré }i}{M_i}\arrow[r]&\underset{\text{et }i-1}{M_i}\arrow[r]&0\arrow[r]&\cdots
\end{tikzcd}\right)
\times\prod_{j\geq0}\left(
\begin{tikzcd}[column sep=small]
0\arrow[r]&\underset{\text{en degré }j}{N_j}\arrow[r]&0\arrow[r]&\cdots
\end{tikzcd}\right)}

De plus, les $M_i$ et $N_j$ de b) (resp. b'))
sont déterminés à isom près par \deleted{\unsure{$K$}} $L_*$,
et ce sont les mêmes (à isom près) dans
b) et b').

\uline{Dém} $\mathrm{a'})\Longleftrightarrow\mathrm{b'})$ est trivial et bien connu,
ainsi que l'unicité à isom \inserted{can.} près des
$M_i$, $N_j$ dans b') (\uline{NB} on a $N_j\simeq H_j(L_*)$
et $M_i\simeq L_i^!/Z_i\ (\simeq B_{i-1})$\unsure{*}.
Or $\mathrm a)\Longleftrightarrow\mathrm{a'})$
\inserted{et $\mathrm b)\Longleftrightarrow\mathrm{b'})$}
par la théorie de Dold-Puppe. D'où
l'équivalence des 4 conditions, et la dernière
assertion est \unsure{contenue|connue}.
